\begin{table}
\caption{Coverage (\%) of 90\% Curie-temperature intervals in every subgroup evaluated}
\label{tab:tableS11_curie_subgroups}
\tablefont
\begin{tabularx}{\textwidth}{@{}>{\raggedright\arraybackslash}Xrrrr@{}}
\toprule
\textbf{Subgroup} & \textbf{Median formulas} & \textbf{Split} & \textbf{Normalized} & \textbf{CQR} \\
\midrule
All test materials & 2,103 & 89.3 & 90.3 & 89.4 \\
Farthest distance decile & 232 & 86.6 & 88.2 & 89.7 \\
Family: Actinide & 47 & 96.2 & 88.9 & 89.6 \\
Family: Intermetallic & 1,137 & 92.9 & 91.4 & 92.9 \\
Family: Oxide & 653 & 92.3 & 89.9 & 91.5 \\
Family: Physics-complete & 81 & 70.2 & 87.0 & 74.5 \\
Family: Rare Earth & 1,117 & 97.0 & 91.8 & 95.5 \\
Family: Transition Metal & 2,020 & 88.9 & 90.1 & 89.2 \\
Hardest sigma decile (defined by the normalized method itself) & 84 & 54.2 & 88.5 & 61.0 \\
Source: text-mined database & 544 & 78.2 & 86.0 & 79.3 \\
Source: curated compilation (anisotropy) & 461 & 92.7 & 92.4 & 92.4 \\
Source: curated compilation (magnetic materials) & 1,534 & 95.4 & 92.2 & 94.9 \\
Top decile of predicted value & 117 & 65.5 & 87.4 & 68.6 \\
\bottomrule
\end{tabularx}
\par\vspace{4pt}\raggedright\noindent Mean over 10 prespecified splits. Constant-width methods over-cover several families (rare earth, actinide, intermetallic) because their intervals are wide everywhere; normalized conformal stays within about 2 points of nominal except for text-mined labels and physics-complete materials.
\end{table}
